% ============================================================
%  Modèle de départ Overleaf — ARS v3.13.0
%  Langue de l'article : English (per user preferences)
% ============================================================
\documentclass[12pt]{article}

% --- Encoding & language --------------------------------------
\usepackage[utf8]{inputenc}     % UTF-8 source encoding
\usepackage[T1]{fontenc}
\usepackage[english]{babel}      % article language = English

% --- Layout: 2.5cm margins, double spacing --------------------
\usepackage[margin=2.5cm]{geometry}
\usepackage{setspace}
\doublespacing

% --- Required packages ----------------------------------------
\usepackage{graphicx}   % figures
\usepackage{float}      % [H] float placement
\usepackage{booktabs}   % publication-quality tables
\usepackage{hyperref}   % clickable links / refs
\usepackage[round]{natbib} % author-year citations

\title{Your Paper Title Here}
\author{Author Name\\ \small Affiliation, email@institution}
\date{\today}

\begin{document}
\maketitle

% ============================================================
%  IMRaD structure (Introduction, Methods, Results, Discussion)
% ============================================================

\begin{abstract}
% A 150–250 word summary: context, gap, method, key result, contribution.
\end{abstract}

\section{Introduction}
% Background, problem statement, research question, and contribution.
% Example citation: prior work shows X \citep{smith2023}.

\section{Methods}
% Data, materials, procedure, and analysis — enough detail to reproduce.

\section{Results}
% Findings only (no interpretation). Reference figures and tables.

% --- Example: inserting a figure from ../../figures/ ----------
% IMPORTANT: reference the PNG version (300 dpi) for portability.
\begin{figure}[H]
    \centering
    \includegraphics[width=0.8\linewidth]{../../figures/png/example_figure.png}
    \caption{Short, self-contained caption describing the figure.}
    \label{fig:example}
\end{figure}
% (A vector .svg twin lives in ../../figures/svg/ for publication.)

% --- Example: a simple booktabs table ------------------------
\begin{table}[H]
    \centering
    \caption{Example table.}
    \label{tab:example}
    \begin{tabular}{lcc}
        \toprule
        Condition & Mean & SD \\
        \midrule
        A & 0.00 & 0.00 \\
        B & 0.00 & 0.00 \\
        \bottomrule
    \end{tabular}
\end{table}

\section{Discussion}
% Interpretation, limitations, implications, and future work.

\section{Conclusion}
% One paragraph: what was learned and why it matters.

% ============================================================
%  Bibliography
% ============================================================
\bibliographystyle{apalike}
\bibliography{exemple}

\end{document}
